
% !TeX root = main.tex
\section{Dynamics Conventions}

\subsection{Purpose}

\texttt{WeightedDynamics} is the collision dynamics implementation documented
by this project.  It is a readable Python implementation whose source-complete
structure is intended to transfer to GLSL with minimal translation.

\subsection{Source-Owned Execution}

Each source particle reads the immutable frame-start position and velocity
snapshots, processes all of its current contacts, vector-sums the planned
contact impulses, and writes only its own velocity and contact state.  A source
must not write a target particle.

The opposing source evaluation reads the same frame-start geometry.  Pair
compatibility and whole-system conservation must therefore emerge from the
contact rule used independently by both sources.

\subsection{Overlap Area Convention}

Overlap area is a geometric contact property.  For a particle-particle contact,
the source-side and target-side overlap areas are the same geometric area, even
when the two particles have different velocities.  Therefore, once the overlap
area is computed for a source-target pair, the target particle's overlap area is
known from the same calculation.

Velocity may affect momentum transfer, phase, or future position, but it does
not make the instantaneous geometric overlap area different for the two
particles in the same contact.

\subsection{Frame State}

Current contact geometry and relative motion are calculated from the immutable
frame-start snapshot:

\begin{itemize}
    \item current positions;
    \item current velocities;
    \item current radii and masses;
    \item current overlap geometry;
    \item the current source-owned contact list.
\end{itemize}

\subsection{Contact Response}

For each source contact, \texttt{WeightedDynamics} calculates relative normal
velocity, reduced-mass closing momentum, overlap area, and a source-local
overlap-area weight.  The current raw impulse is:

\begin{equation}
J_{\mathrm{raw},k}
=
q_k A_k
\left(\frac{A_k}{\sum_j A_j}\right)
\Delta t.
\end{equation}

Compression is capped by the contact's weighted share of source available
closing momentum.  Release is capped by the contact-owned stored internal
momentum.  All planned contact impulses are vector-summed before the source
velocity is written.

This equation documents the current implementation.  It does not assert that
source-local weighting conserves whole-system vector momentum or kinetic
energy; those properties must be established by diagnostics and test results.

Current \texttt{ThreeParticle} diagnostics show that the weighted rule is not
conservative.  The weighted capture finishes with approximately
\(0.04899135\) vector momentum drift and \(+0.00090284\) kinetic-energy drift,
even though total stored internal momentum returns to zero.  Source-local area
weighting is therefore a known model failure, not an accepted conservation
rule.

\subsection{Persistent Contact State}

\texttt{WeightedDynamics} currently stores source-owned internal momentum and
collision phase by target ID.  Current-frame contact slots are rebuilt each
frame; the target-keyed dictionaries retain the persistent contact state used
for compression and release.

\subsection{Diagnostics}

Momentum and kinetic energy are reported diagnostics.  They are used to
evaluate the contact rule but are not direct dynamics inputs.  Documentation
must distinguish actual vector momentum drift from scalar internal-momentum
ledger totals.  A drained internal-momentum ledger does not, by itself, prove
momentum or kinetic-energy conservation.
